% Generated draft: requires mathematical and source review.
\documentclass[11pt]{article}
\usepackage[margin=25mm]{geometry}
\usepackage{fontspec}
\setmainfont{Latin Modern Roman}
\usepackage{amsmath,amssymb,mathtools,unicode-math}
\setmathfont{Latin Modern Math}
\usepackage{enumitem,needspace,xurl,hyperref}
\setlength{\parindent}{0pt}
\setlength{\parskip}{5pt}
\setlength{\emergencystretch}{3em}
\newcommand{\N}{\mathbb{N}}
\newcommand{\Z}{\mathbb{Z}}
\newcommand{\Q}{\mathbb{Q}}
\newcommand{\R}{\mathbb{R}}
\newcommand{\C}{\mathbb{C}}
\newcommand{\F}{\mathbb{F}}
\newcommand{\A}{\mathbb{A}}
\newcommand{\PP}{\mathbb{P}}
\newcommand{\E}{\mathbb{E}}
\newcommand{\eps}{\varepsilon}
\newcommand{\dd}{\,\mathrm{d}}
\begin{document}
\section*{UnsolvedMath ID: 700157}
\textbf{Source problem code:} catalogue entry
\subsection{Definitions and mathematical statement}
The supplied problem does not specify the coefficient field or ring, its characteristic, whether the algebra has an identity, or whether automorphisms must fix the coefficients. Thus the following formulation is conditional on a standard interpretation, rather than a resolution of these omissions.

Fix a field \(K\). For an integer \(n\geq 2\), let
\[
A_n(K)=K\langle x_1,\ldots,x_n\rangle
\]
be the free unital associative \(K\)-algebra: its elements are finite \(K\)-linear combinations of words in the generators, including the empty word, and multiplication is induced by concatenation. Write \(\operatorname{Aut}_{K}(A_n(K))\) for the group, under composition, of bijective unital \(K\)-algebra homomorphisms from \(A_n(K)\) to itself.

\textbf{Problem (L. A. Bokut', 1969; Kourovka Notebook 3.3).} Under this interpretation, describe
\[
\operatorname{Aut}_{K}\bigl(K\langle x_1,\ldots,x_n\rangle\bigr)
\qquad\text{for each integer }n\geq 2.
\]
The supplied source calls this a well-known problem. It does not specify the intended form of a description, nor whether the coefficient field is fixed or allowed to vary. These qualifications remain unresolved.

For interpreting the supplied background under the same convention, an affine automorphism sends
\[
x_i\longmapsto \sum_{j=1}^{n}a_{ij}x_j+b_i,
\qquad (a_{ij})\in\operatorname{GL}_n(K),\quad b_i\in K.
\]
An elementary automorphism fixes every generator other than some \(x_i\) and sends
\[
x_i\longmapsto a x_i+f,
\qquad a\in K\setminus\{0\},\quad
f\in K\langle x_1,\ldots,x_{i-1},x_{i+1},\ldots,x_n\rangle.
\]
The subgroup generated by affine and elementary automorphisms is called the tame subgroup; an automorphism outside it is called wild.

The supplied record reports that all automorphisms are tame when \(n=2\), attributing this to A. J. Czerniakiewicz and L. G. Makar-Limanov. It also reports that the Anick automorphism is wild when \(n=3\), attributing this to U. U. Umirbaev. The coefficient and characteristic hypotheses of these reports are not supplied, and no formula for the Anick automorphism is given. Finally, it reports an algorithm, attributed to A. V. Jagzhev, for recognizing which endomorphisms of free associative algebras of finite rank over constructive fields are automorphisms. The effective presentation meant by \emph{constructive fields} is not specified. These are supplied background assertions, not independently verified results or additional requests to solve the problem.
\subsection{Short English statement}
Describe the automorphism group of a free associative algebra on \(n\geq 2\) generators. A conditional formulation uses unital algebras over a fixed field \(K\) and automorphisms fixing \(K\); the source leaves these conventions and the intended form of the description unspecified. Its supplied background reports tameness in rank two, a wild Anick automorphism in rank three, and an automorphism-recognition algorithm over constructive fields.
\subsection{Sources}
\begin{itemize}
\item Local UnsolvedMath registry: permanent ID \(700157\), exact source problem code \texttt{catalogue entry}, title \emph{Kourovka Notebook Problem 3.3}. The supplied registry labels the problem \emph{open}; this label is supplied curation, not an independently verified assessment.
\item \emph{The Kourovka Notebook (Unsolved Problems in Group Theory)}, original problem 3.3, third issue (1969), as identified in the supplied record with the twenty-first issue edition (2026), \texttt{arXiv:1401.0300v48}. Proposed by L. A. Bokut'. The supplied collection description names E. I. Khukhro and V. D. Mazurov as editors of Issues 1--21 (1965--2026).
\item A. J. Czerniakiewicz, \emph{Trans. Amer. Math. Soc.} \textbf{160} (1971), 393--401; \textbf{171} (1972), 309--315. Supplied references for tameness in rank two; article titles and hypotheses are not supplied.
\item L. G. Makar-Limanov, \emph{Funct. Anal. Appl.} \textbf{4} (1971), 262--264. Supplied reference for tameness in rank two; the article title and hypotheses are not supplied.
\item U. U. Umirbaev, \emph{J. Reine Angew. Math.} \textbf{605} (2007), 165--178. Supplied reference for wildness of the Anick automorphism in rank three; the article title and hypotheses are not supplied.
\item A. V. Jagzhev, \emph{Siberian. Math. J.} \textbf{21} (1980), 142--146. Supplied reference for recognizing automorphisms among endomorphisms of finite-rank free associative algebras over constructive fields; the article title and precise effectiveness assumptions are not supplied.
\item The registry records creation and updating at \texttt{2026-10-16T14:00:00Z} and upstream revision \texttt{d371ee917ac9647fbb06cd6fc529deec6f3dda82}. These are registry provenance metadata, not a literature-verification date. No \texttt{literature\_checked\_at} value or external reference URLs are supplied. The references above have not been independently inspected for this draft.
\end{itemize}
\end{document}
